\documentclass[10pt,a4paper]{amsart}
\usepackage[margin=19mm]{geometry}
\usepackage{amsmath,amssymb,mathtools}
\usepackage[scr=rsfs,cal=euler]{mathalfa}
\usepackage{xfrac}
\usepackage[unicode,hidelinks,bookmarks=false]{hyperref}
\newcommand{\ie}{i.e.\ }
\newcommand{\cf}{cf.\ }
\newcommand{\resp}{respectively\ }
\newcommand{\Subsection}{\subsection}
\providecommand{\nobreakdash}{\nobreak\hbox{-}}
\setlength{\emergencystretch}{2em}
\multlinegap0pt
\usepackage{bm}
\usepackage{bbm}
\usepackage{dsfont}
\usepackage{xspace}
\usepackage{cancel}
\usepackage{mathtools}
\usepackage{amsthm}
\makeatletter
\@ifundefined{equation}{\newtheorem{equation}{Equation}}{}
\@ifundefined{Lem}{\newtheorem{Lem}[equation]{Lemma}}{}
\makeatother
\newcommand{\End}{\operatorname{End}}
\newcommand{\Gt}{\widetilde{G}}
\newcommand{\Hcal}{\mathcal{H}}
\newcommand{\Hom}{\operatorname{Hom}}
\newcommand{\Ift}{\widetilde{\mathsf{I}}}
\newcommand{\Irr}{\operatorname{Irr}}
\newcommand{\It}{\widetilde{I}}
\newcommand{\Tt}{\widetilde{T}}
\newcommand{\Uf}{\mathsf{U}}
\title[Minimal depth $K$-types for wild double covers and Shimura c]{Minimal depth $K$-types for wild double covers and Shimura correspondences}
\author{Edmund Karasiewicz, Shuichiro Takeda}
\date{Source: arXiv:2510.23265v2 (CC BY 4.0)}
\begin{document}
\maketitle

\noindent\emph{Source and licence.} Minimal depth $K$-types for wild double covers and Shimura correspondences. Version arXiv:2510.23265v2 (CC BY 4.0), distributed under the Creative Commons Attribution 4.0 International licence, as recorded in the arXiv licence metadata for this exact version and shown on the abstract page https://arxiv.org/abs/2510.23265v2. Full rights notice: licenses/arxiv-2510.23265-CC-BY-4.0.txt.\par\smallskip

\noindent\textit{Editorial scope.} Counted unit: the Lem whose source label is \texttt{LemSurj} in the article (source0.tex lines 4444--4446, source label \texttt{LemSurj}), delivered together with the article's own complete original proof of it (source0.tex lines 4448--4527, one \texttt{proof} environment of 80 lines). No mathematics is written, repaired or re-derived: statement and proof are the article's own text line for line. Required context retained uncounted: L3to1 (lines 4414--4416); VeryPosSupp (lines 4424--4426). Every definition, display and label of those lines is retained unchanged. Omitted from this package and not redistributed: 1--4413, 4417--4423, 4427--4443, 4447--4447, 4528--4915. Nothing inside a retained excerpt is omitted or altered: every retained body is delivered exactly as the article's lines are written, so no omission receipt of any kind accompanies this package. Citations: the retained excerpts carry no citation command at all, so no bibliography entry of the article is transcribed and no reference is redistributed. Reference adaptation: none was needed and none was made; every reference inside the delivered unit resolves inside it, so no unresolved reference or citation is produced. Printed slips kept literal: no printed word, symbol, hypothesis or number of the counted statement or of its proof was altered. Layout and class: the article's own class is \texttt{amsart} with packages including amscd, amssymb, enumerate, xypic, multirow, hhline, verbatim, mathdots, comment, tikz-cd, appendix, stmaryrd, graphicx, mathtools; the wrapper is the lane's pinned \texttt{amsart} editorial wrapper at 10pt on A4, which restates the theorem-like environments the retained text needs and adds the article's own macro definitions that the retained text needs (\texttt{\textbackslash End}, \texttt{\textbackslash Gt}, \texttt{\textbackslash Hcal}, \texttt{\textbackslash Hom}, \texttt{\textbackslash Ift}, \texttt{\textbackslash Irr}, \texttt{\textbackslash It}, \texttt{\textbackslash Tt}, \texttt{\textbackslash Uf}), with extra wrapper packages (bm, bbm, dsfont, xspace, cancel, mathtools). The delivered numbering is the unit's own. Assets: no retained span contains a figure environment, a graphics-inclusion command, a PDF-page inclusion, a table or a TikZ picture; no image file is copied into this package, the package declares no asset dependency and no image file is redistributed with it. Licence: version arXiv:2510.23265v2 of this article is distributed under the Creative Commons Attribution 4.0 International licence, as recorded in the arXiv licence metadata for this exact version and shown on the abstract page https://arxiv.org/abs/2510.23265v2. A full rights notice is delivered at licenses/arxiv-2510.23265-CC-BY-4.0.txt. Characters: every retained excerpt is pure ASCII, so the delivered engine, XeLaTeX with its default Latin Modern text font, reports no missing character.
\begin{Lem}\label{L3to1}
	Let $G$ be a finite group with an Iwahori factorization $(V,T,U)$. Let $\pi\in \Irr(G)$ be such that $r_{U,V}\pi\neq 0$. The element $e_{U}e_{V}e_{U}\in \End({\pi})$ induces a linear isomorphism of the space $\pi^{U}$ and thus is a generator of the one-dimensional space $\Hom_{T}(\pi^{U},\pi^{U})=\mathbb{C}e_{U}$.
\end{Lem}

\begin{Lem}\label{VeryPosSupp}
	Let $z\in Z(\Tt)$ be such that $zU_{0}z^{-1}\subseteq U_{2e}$. Then the element $e_{\Uf_{0}}e_{\Uf^{-}_{1}}\in\Hcal(\,\Ift\,)$ defines a nonzero element in $\Hom_{\It\cap\,^{z}\It}(\sigma,\,^{z}\sigma)$.
\end{Lem}

\begin{Lem}\label{LemSurj}
Let $q:\pi\to\pi_U$ be the canonical surjection. The map $\Hom(\sigma,\pi)\rightarrow \Hom(\sigma^{U_{0}},\pi_{U})$ defined by $f\mapsto q\circ f|_{\sigma^{U_{0}}}$ is a homomorphism of $(\It\cap \widetilde{B},\widetilde{B})$-bimodules. The map induced by restriction $\Psi:\Hom_{\It}(\sigma,\pi)\rightarrow \Hom_{\Tt_{0}}(\sigma^{U_{0}},\pi_{U})$ is a surjective linear map.
\end{Lem}

\begin{proof}
	Let $f\in \Hom_{\Tt_{0}}(\sigma^{U_{0}},\pi_{U})$. Note that by definition $f\circ e_{{\Uf}_{0}}=f$. When convenient we use this to identify $f$ with $f\circ e_{{\Uf}_{0}}\in \Hom_{\Tt_{0}}(\sigma,\pi_{U})$.
	
	Now since the category of $\Tt_{0}$-modules is semisimple we can fix a $\Tt_{0}$-module splitting $s:\pi_{U}\rightarrow \pi$ of the $\Tt_{0}$-module surjection $q:\pi\rightarrow \pi_{U}$. Define 
    \[
    F:=s\circ f\circ e_{{\Uf}_{0}}\in \Hom_{\Tt_{0}}(\sigma,\pi),
    \]
    so that $q\circ F=f\circ e_{{\Uf}_{0}}$.
	
	Since $\sigma$ is finite dimensional, there exists $j\in \mathbb{Z}_{\geq 1}$ such that $\Im(F)\subseteq \pi^{U^{-}_{j}}$. Let $z\in Z(\Tt)$ be such that $z^{-1}U^{-}_{1}z\subseteq U^{-}_{j}\cap U^{-}_{2e+1}$.

    Define
    \[
    F_{1}:=\int_{\It}\pi(x)\pi(z)F\circ \sigma(x^{-1})dx\in\Hom_{\It}(\sigma,\pi).
    \]
    Let us compare $F_1$ to $f$. Let $w\in \sigma$. Then by the Iwahori factorization, we have
	\begin{align*}
		q\circ F_1(e_{{\Uf}_{0}}w)=&q\circ \int_{\It} \pi(x)\pi(z)F\circ \sigma(x^{-1}) e_{{\Uf}_{0}}w\,dx\\
		=&q\circ \int_{U_{0}\Tt_{0}U^{-}_{1}} \pi(utv)\pi(z)F\circ \sigma((utv)^{-1}) e_{{\Uf}_{0}}w \,du\,dt\,dv\\
		=&q\circ \int_{U_{0}\Tt_{0}U^{-}_{1}} \pi(utvz)F\circ \sigma((utv)^{-1}) e_{{\Uf}_{0}}w \,du\,dt\,dv.
	\end{align*}
	Using the idempotent $e_{{\Uf}_{0}}$ and the fact that $F$ is a $\Tt_{0}$-module homomorphism we have 
	\begin{align*}
		&q\circ \int_{U_{0}\Tt_{0}U^{-}_{1}} \pi(utvz)F\circ \sigma((utv)^{-1}) e_{U_{0}}w \,du\, dt\, dv\\
        \equiv&q\circ \int_{\Tt_{0}U^{-}_{1}} \pi(tvz)F\circ \sigma((tv)^{-1}) e_{{\Uf}_{0}}w \,dt\,dv\\
		\equiv&q\circ \int_{U^{-}_{1}} \pi(vz)F\circ \sigma(v^{-1}) e_{{\Uf}_{0}}w \,dv,
	\end{align*}
    where by $\equiv$ we mean equality up to a positive scalar. (The discrepancy comes from the choice of measures. We will use the same convention in what follow.)
	Now because $\Im(F)\subseteq \pi^{U^{-}_{j}}$ and  $z^{-1}U^{-}_{1}z\subseteq U^{-}_{j}$, we have
	\begin{align*}
		q\circ \int_{U^{-}_{1}} \pi(vz)F\circ \sigma(v^{-1}) e_{{\Uf}_{0}}w \,dv
        \equiv&q\circ \int_{U^{-}_{1}} \pi(z)\pi(z^{-1}vz)F\circ \sigma(v^{-1}) e_{{\Uf}_{0}}w \,dv\\
		\equiv&q\circ \pi(z)F\circ e_{{\Uf}^{-}_{1}}e_{{\Uf}_{0}}w\\
		\equiv&\pi_{U}(z)f \circ e_{{\Uf}_{0}}e_{{\Uf}^{-}_{1}}e_{{\Uf}_{0}}w,
	\end{align*}
    where the last equality is by the definition of $F$. Finally we can apply Lemma \ref{L3to1} to get 
	\begin{equation*}
		\pi_{U}(z)f \circ e_{{\Uf}_{0}}e_{{\Uf}^{-}_{1}}e_{{\Uf}_{0}}w\equiv\pi_{U}(z)f \circ e_{{\Uf}_{0}}w.
	\end{equation*}
	Thus we have shown that
	\begin{equation}\label{SurjEq1}
		q\circ F_1(e_{{\Uf}_{0}}w)\equiv\pi_{U}(z)f \circ e_{{\Uf}_{0}}w.
	\end{equation}
	
	Recall that for $T\in \Hcal$ and $A\in \Hom_{\It}(\sigma,\pi)$, we define $A*T\in\Hom_{\It}(\sigma,\pi)$ by
    \[
    (A*T)(w):=\int_{\Gt}\pi(h^{-1})A(T(h)w)\,dh.
    \]
    Now suppose that $T\in \Hcal_{z^{-1}}$ is such that $T(z^{-1})=e_{{\Uf}_{0}}e_{{\Uf}^{-}_{1}}$. Such an element exists by Lemma \ref{VeryPosSupp} because $z^{-1}U^{-}_{1}z\subseteq U^{-}_{2e+1}$. Then by definition 
	\begin{align*}
		q[A*T(w)] =&q[ \int_{\Gt}\pi(h^{-1})A(T(h)w)\,dh]\\
		=& q[\int_{\It z^{-1}\It}\pi(h^{-1})A(T(h)w)\,dh]\\
		=& q[\sum_{\It\backslash \It z^{-1}\It}\pi(h^{-1})A(T(h)w)].
	\end{align*}
	Now because $z^{-1}U^{-}_{1}z\subseteq U^{-}_{j}\cap U^{-}_{2e+1}$,
	\begin{align*}
		q[\sum_{\It\backslash \It z^{-1}\It}\pi(h^{-1})A(T(h)w)] =& q[\sum_{zU_{0}z^{-1}\backslash U_{0}}\pi(u^{-1}z)A(T(z^{-1})\sigma(u)w)]\\
		=& q[\sum_{zU_{0}z^{-1}\backslash U_{0}}\pi(z)A(e_{{\Uf}_{0}}e_{{\Uf}^{-}_{1}}\sigma(u)w)]\\
		\equiv&\pi_{U}(z)q[A(e_{{\Uf}_{0}}e_{{\Uf}^{-}_{1}}e_{{\Uf}_{0}}w)].
	\end{align*}
	By Lemma \ref{L3to1},
	\begin{equation*}
	\pi_{U}(z)q[A(e_{{\Uf}_{0}}e_{{\Uf}^{-}_{1}}e_{{\Uf}_{0}}w)]\equiv\pi(z)q[A(e_{{\Uf}_{0}}w)].
	\end{equation*}
	Thus we have shown
	\begin{equation}\label{SurjEq2}
		q[A*T(w)] \equiv\pi_{U}(z)q[A(e_{{\Uf}_{0}}w)].
	\end{equation}
	
	Now we set $A=F_1*T^{-1}$ and combine equations \eqref{SurjEq1} and \eqref{SurjEq2} to get 
    \begin{equation*}
		\pi_{U}(z)f \circ e_{{\Uf}_{0}}w\stackrel{\eqref{SurjEq1}}{\equiv}q\circ F_1(e_{{\Uf}_{0}}w)\equiv q\circ F_1*T^{-1}*T(e_{{\Uf}_{0}}w)\stackrel{\eqref{SurjEq2}}{\equiv}\pi_{U}(z)q[[F*T^{-1}](e_{{\Uf}_{0}}w)].
	\end{equation*}
	
	Thus 
	\begin{equation*}
		q\circ [F*T^{-1}]\circ e_{{\Uf}_{0}}\equiv f\circ e_{{\Uf}_{0}}
	\end{equation*}
	and the result follows.
\end{proof}

\end{document}
